\chapter{THE FOUR REQUIRED ALGORITHMS}\nonum

\section{Bresenham Line Drawing}
A pixel grid has whole-number positions, but a real line passes between them. For each step along the line we must pick the pixel nearest to the line \cite{bresenham1965}. Call the longer direction the major axis and the shorter one the minor axis, with lengths $\Delta_M$ and $\Delta_m$. We move one pixel at a time along the major axis. At each step the minor coordinate either stays the same (a straight step) or changes by one (a diagonal step). A whole-number error value decides which:
\[
e_0 = 2\Delta_m-\Delta_M,\qquad
e_{i+1}=\begin{cases} e_i+2(\Delta_m-\Delta_M) & e_i>0\ \ (\text{diagonal})\\ e_i+2\Delta_m & e_i\le 0\ \ (\text{straight}).\end{cases}
\]
Only whole-number additions and comparisons are needed, with no division and no rounding. When $e=0$ the line is exactly half way between two pixels. We then pick the larger minor coordinate, so a line drawn backwards gives the same pixels as one drawn forwards.

\begin{figure}[H]
\centering
\begin{lstlisting}[style=code]
const diagonal = error > 0 || (error === 0 && minorSign > 0)
if (diagonal) { error += 2 * (minor - major); x += sx; y += sy }
else          { error += 2 * minor; if (steep) y += sy; else x += sx }
\end{lstlisting}
\caption{The decision in \texttt{bresenhamSteps} (\texttt{src/core/render/line.ts})}
\label{fig:brescode}
\end{figure}

\textbf{Worked example.} For the line from $(0,0)$ to $(5,2)$ we have $\Delta_M=5$, $\Delta_m=2$ and $e_0=-1$ (Table~\ref{tab:bres}). The unit test checks exactly these values. We also kept the older DDA method in the code, only so the lab can draw it over Bresenham's line (Figure~\ref{fig:bres}).

\begin{table}[H]
\centering
\caption{Bresenham steps for the line $(0,0)\to(5,2)$}
\label{tab:bres}
\small
\begin{tabular}{@{}cccc@{}}
\toprule
Pixel & Error before & Move & Error after\\
\midrule
(0,0) & $-1$ & straight & 3\\
(1,0) & 3 & diagonal & $-3\ $\\
(2,1) & $-3$ & straight & 1\\
(3,1) & 1 & diagonal & $-5\ $\\
(4,2) & $-5$ & straight & $-1\ $\\
(5,2) & $-1$ & none & $-1\ $\\
\bottomrule
\end{tabular}
\end{table}

\twoshot{02-bresenham-step-table.png}{02b-bresenham-dda-overlay.png}{Bresenham Lab. Left: a line half drawn with its table of error values. Right: the DDA line drawn over it.}{fig:bres}

\section{Back-Face Culling}\label{sec:cull}
On a closed solid, a triangle that faces away from the camera is always hidden behind the triangles that face it, so we can skip it. The front of a triangle is the side from which its corners $a\to b\to c$ go anticlockwise. With $e_1=b-a$ and $e_2=c-a$, the normal is $\mathbf n=e_1\times e_2$. Its length is twice the triangle's area, which is the length of the outer product $e_1\wedge e_2$. The vector from the eye $E$ to the triangle is $a-E$, and the triangle is back-facing when
\[
\mathbf n\cdot(a-E)\ \ge\ 0 .
\]
The lab also has a second test that uses the projected corners. The doubled signed area $(x_1-x_0)(y_2-y_0)-(y_1-y_0)(x_2-x_0)$ is negative for a front-facing triangle, because screen $y$ points down. Both tests give the same answer. A triangle seen exactly edge-on gives 0 in exact maths, but rounding can give a tiny value like $10^{-17}$ of either sign. Both tests treat anything closer to zero than $10^{-9}$ as back-facing, so they always agree.

\textbf{Worked example.} Let $a=(0,0,0)$, $b=(1,0,0)$, $c=(0,1,0)$ and the eye $E=(0,0,5)$. Then $\mathbf n=(0,0,1)$, the area is $0.5$, and $\mathbf n\cdot(a-E)=-5$. This is below zero, so the triangle is front-facing. If we list the corners as $a,c,b$, the normal becomes $(0,0,-1)$ and the product is $+5$, so the triangle is back-facing. Our tests check both orders. About half of the triangles of a closed shape face away (Figure~\ref{fig:cull}).

\twoshot{03-culling-on-wireframe-normals.png}{03b-culling-off-wireframe.png}{Culling Lab, wireframe with normals. Left: culling on. Right: culling off, so the back faces are drawn too.}{fig:cull}

\section{Alpha Compositing}
A colour with alpha has red, green and blue from 0 to 255 and an alpha from 0 (invisible) to 1 (solid). The ``over'' rule says what we see when a see-through colour (the source) is laid over another colour (the destination) \cite{porterduff1984}. The source hides a share $\alpha_s$ of what is below, so a share $1-\alpha_s$ of the destination still shows:
\[
\alpha_o=\alpha_s+\alpha_d(1-\alpha_s),\qquad
C_o=\frac{C_s\alpha_s+C_d\alpha_d(1-\alpha_s)}{\alpha_o}.
\]
We divide by $\alpha_o$ to turn the weighted sum back into a normal colour. If $\alpha_o=0$ we return transparent black, to avoid dividing by zero. The code also has a premultiplied version, where every channel uses one rule: $X_o=X_s+X_d(1-\alpha_s)$. A test checks that both versions give the same picture.

\begin{figure}[H]
\centering
\begin{lstlisting}[style=code]
const outA = as + ad * (1 - as)
if (outA === 0) return [0, 0, 0, 0]
const dstWeight = ad * (1 - as)
return [(src[0]*as + dst[0]*dstWeight) / outA, /* same for g and b */ ..., outA]
\end{lstlisting}
\caption{The core of \texttt{over} (\texttt{src/core/render/alpha.ts})}
\label{fig:overcode}
\end{figure}

\textbf{Worked example.} A half-transparent red $(255,0,0,0.5)$ over a solid blue $(0,0,255,1)$ gives $\alpha_o=1$ and the colour $(127.5,0,127.5)$, a purple. A solid blue over the red gives pure blue, because a solid source hides everything. So the order of the layers matters.

For many see-through triangles in 3D, the renderer draws in two passes. Solid triangles come first and write the depth buffer. See-through triangles come second. They use the depth test (so a pane behind a solid object is hidden) but do not write depth (so panes do not hide each other). They are sorted from far to near, so the farthest is blended first. The lab can switch the sort off, which gives wrong colours (Figure~\ref{fig:alpha3d}).

\twoshot{04b-alpha-3d-sort-on.png}{04c-alpha-3d-sort-off.png}{Alpha Lab in 3D. Left: panes sorted back to front. Right: sorting off, so the colours are wrong.}{fig:alpha3d}

\section{Visvalingam--Whyatt Simplification}
A polyline with many points has points that hardly bend it. For each point we look at the triangle it makes with its two neighbours. A small triangle means the point is not needed. We remove the point with the smallest triangle, recompute the triangles of its two neighbours, and repeat until enough points are left \cite{visvalingam1993}. The area of a triangle is half the 2D outer product of two of its edges.

The two end points of an open line are never removed. A closed polygon keeps at least three points. If a neighbour's new triangle is smaller than the one just removed, we rank it with the removed area instead: ranking area $=\max(\text{new area},\ \text{last removed area})$. This keeps the removal areas from ever going down, so ``remove all points below area $X$'' is the same as ``remove the first $k$ points''. Ties go to the lower index.

Our version scans all points for the smallest area after every removal. This costs $O(n^2)$ for $n$ points. A heap would give $O(n\log n)$, but the scan is short and easy to explain, and Chapter~7 measures its cost.

\begin{figure}[H]
\centering
\begin{lstlisting}[style=code]
for (let i = 0; i < n; i++)
  if (alive[i] && canGo(i) && (best < 0 || area[i] < area[best])) best = i
...
area[v] = Math.max(rawArea[v], lastArea)   // the rule that areas never go down
\end{lstlisting}
\caption{The scan and the ranking rule (\texttt{src/core/simplify/visvalingam.ts})}
\label{fig:vwcode}
\end{figure}

\textbf{Worked example.} Take the open line $(0,0),(1,1),(2,0),(3,3)$. Point 1 makes the triangle $(0,0),(1,1),(2,0)$ with area 1. Point 2 makes $(1,1),(2,0),(3,3)$ with area 2. Point 1 goes first, with area 1. Point 2 now has the neighbours $(0,0)$ and $(3,3)$, a triangle of area 3, and it is ranked $\max(3,1)=3$. Our function returns these two removals with areas 1 and 3. We also kept Ramer--Douglas--Peucker \cite{ramer1972,douglas1973} for comparison (Figure~\ref{fig:vw}).

\twoshot{05-simplify-500-vs-75-rdp.png}{05b-simplify-step-mode.png}{Simplify Lab. Left: 500 points against 75, with the RDP result drawn over it. Right: step mode, one removal at a time.}{fig:vw}
